% !TeX program = xelatex
% 本模板用于全国大学生数学建模竞赛 AI 工具使用支撑材料。
% 仅填写真实、经队员审核的内容；按实际条目增删表格行和记录块。
\documentclass[UTF8,a4paper,11pt]{ctexart}
\usepackage[a4paper,top=2.2cm,bottom=2.2cm,left=2.15cm,right=2.15cm]{geometry}
\usepackage{array,booktabs,longtable,tabularx,fancyhdr,lastpage,enumitem}
\usepackage[table]{xcolor}
\usepackage[hidelinks]{hyperref}

\definecolor{rulegray}{HTML}{B7BDC5}
\definecolor{headgray}{HTML}{EEF1F4}
\definecolor{placeholder}{HTML}{666666}
\newcommand{\fillin}[1]{\textcolor{placeholder}{【#1】}}
\newcolumntype{Y}{>{\raggedright\arraybackslash}X}
\newcolumntype{L}[1]{>{\raggedright\arraybackslash}p{#1}}
\renewcommand{\arraystretch}{1.35}
\setlength{\parindent}{2em}
\setlength{\parskip}{0.25em}
\setlist[itemize]{leftmargin=2em,itemsep=0.2em,topsep=0.3em}

\pagestyle{fancy}
\setlength{\headheight}{15pt}
\fancyhf{}
\fancyhead[L]{全国大学生数学建模竞赛支撑材料}
\fancyhead[R]{AI 工具使用详情}
\fancyfoot[C]{第 \thepage\ 页，共 \pageref{LastPage} 页}
\renewcommand{\headrulewidth}{0.4pt}
\renewcommand{\footrulewidth}{0pt}

\begin{document}

\begin{center}
  {\zihao{2}\bfseries AI 工具使用详情}\par
  \vspace{0.4em}
  {\zihao{5}\color{placeholder}填写模板（提交前删除本行）}
\end{center}

\section*{一、基本信息}
\noindent
\begin{tabularx}{\textwidth}{L{3.6cm}Y}
\toprule
\rowcolor{headgray}\textbf{项目} & \textbf{填写内容} \\
\midrule
竞赛名称 & \fillin{填写竞赛全称} \\
赛题编号或名称 & \fillin{填写赛题编号或名称} \\
参赛队号 & \fillin{填写参赛队号} \\
记录起止时间 & \fillin{年/月/日—年/月/日} \\
材料生成日期 & \fillin{年/月/日} \\
\bottomrule
\end{tabularx}

\section*{二、责任说明与填写原则}
\noindent
\fcolorbox{rulegray}{headgray}{%
\parbox{0.94\textwidth}{
本材料记录本队在竞赛过程中实际使用的 AI 工具及处理方式。核心建模与分析由参赛队主导；AI 生成或建议内容均须由队员人工审查、修改并验证。不得写入 API 密钥、账号密码、完整聊天记录、模型私有思维链或无关个人信息。
}}

\section*{三、AI 工具目录}
\begin{longtable}{@{}L{1.4cm}L{3.0cm}L{3.0cm}L{3.2cm}L{4.0cm}@{}}
\toprule
\rowcolor{headgray}\textbf{编号} & \textbf{工具名称} & \textbf{提供方} & \textbf{版本或模型} & \textbf{使用入口或方式} \\
\midrule
\endfirsthead
\toprule
\rowcolor{headgray}\textbf{编号} & \textbf{工具名称} & \textbf{提供方} & \textbf{版本或模型} & \textbf{使用入口或方式} \\
\midrule
\endhead
TOOL-01 & \fillin{名称} & \fillin{提供方} & \fillin{界面或官方信息所示版本/模型} & \fillin{网页端、客户端、API 等} \\
\bottomrule
\end{longtable}

\section*{四、分阶段使用记录}
\subsection*{AI-001：\fillin{阶段名称}}
\noindent
\begin{tabularx}{\textwidth}{L{4.2cm}Y}
\toprule
\rowcolor{headgray}\textbf{记录项} & \textbf{真实使用情况} \\
\midrule
所用工具编号 & \fillin{如 TOOL-01} \\
使用目的 & \fillin{说明该阶段为何使用 AI} \\
主要提示方式或代表性提示摘要 & \fillin{概括提示结构和关键约束，不复制全部对话} \\
使用过程摘要 & \fillin{说明主要步骤、迭代和人工参与} \\
AI 输出采纳情况 & \fillin{完全采纳、部分采纳或未采纳，并说明范围} \\
人工修改内容 & \fillin{说明队员做出的实质修改} \\
人工核验方法与结果 & \fillin{说明复算、对照、测试、文献核查等方法及结果} \\
对应证据文件或产物 & \fillin{填写可定位的文件、图表、日志或章节} \\
负责确认的队员与状态 & \fillin{队员；已确认} \\
\bottomrule
\end{tabularx}

\vspace{0.6em}
{\small\color{placeholder}如有多个阶段，复制以上记录块并依次编号为 AI-002、AI-003……}

\section*{五、典型交互示例}
\begin{longtable}{L{1.6cm}L{4.2cm}L{4.2cm}L{4.2cm}}
\toprule
\rowcolor{headgray}\textbf{对应记录} & \textbf{代表性提示摘要} & \textbf{AI 结果摘要} & \textbf{队员处理} \\
\midrule
\endfirsthead
\toprule
\rowcolor{headgray}\textbf{对应记录} & \textbf{代表性提示摘要} & \textbf{AI 结果摘要} & \textbf{队员处理} \\
\midrule
\endhead
AI-001 & \fillin{摘要} & \fillin{摘要} & \fillin{采纳、修改与核验} \\
\bottomrule
\end{longtable}

\section*{六、未采纳的重要 AI 建议或结果}
\begin{itemize}
  \item \fillin{说明未采纳内容、原因及人工判断；如无则填写“无”}
\end{itemize}

\section*{七、人工修改与验证汇总}
\noindent
\begin{tabularx}{\textwidth}{L{3.6cm}Y}
\toprule
\rowcolor{headgray}\textbf{项目} & \textbf{说明} \\
\midrule
主要人工修改 & \fillin{汇总队员对 AI 输出所作修改} \\
主要核验方法 & \fillin{汇总复算、代码测试、交叉验证、文献核对等} \\
核验依据或证据 & \fillin{列出关键证据文件和结果} \\
仍有待核实内容 & \fillin{提交前应为“无”} \\
\bottomrule
\end{tabularx}

\section*{八、最终确认}
\noindent
\begin{tabularx}{\textwidth}{L{7.2cm}Y}
\toprule
\rowcolor{headgray}\textbf{确认事项} & \textbf{确认结果} \\
\midrule
核心建模与分析由参赛队主导 & \fillin{是/否} \\
所有记录已逐项人工审查 & \fillin{是/否} \\
仍有待核实内容 & \fillin{是/否；如是请说明} \\
队员共同确认 & \fillin{填写确认方式或队员姓名} \\
确认日期 & \fillin{年/月/日} \\
\bottomrule
\end{tabularx}

\vfill
{\small\color{placeholder}提交前检查：文件名必须为“AI 工具使用详情.pdf”；删除模板提示；确认无占位符、密钥和未经审核的陈述。}

\end{document}
